% Chapter 21: Active and Passive Voice: Who Does It, Who Gets It

\chapter{Active and Passive Voice: Who Does It, Who Gets It}

\section{Pedagogical Purpose}
Having learned to build fuller, more skilful sentences in Chapter 20, learners are ready to see that the same event can be described in two different ways, depending on whether the sentence focuses on the doer or the receiver of the action. This is an early, carefully limited introduction to voice — the goal is recognition and simple transformation, not full mastery of every passive form.

\begin{learningobjectives}
The learner can:
\begin{itemize}
    \item identify the doer and the receiver of an action in a sentence.
    \item recognise whether a sentence is in the active or passive voice.
    \item change a simple active sentence into the passive voice, and back again.
    \item notice when passive voice is used in real life (e.g., notices, news headlines).
\end{itemize}
\end{learningobjectives}

\section{Prerequisite Check}
\begin{quickcheck}
Underline the subject and circle the verb in each sentence.
\begin{enumerate}
    \item \underline{Rohan} \textcircled{kicked} the ball.
    \item \underline{Maya} \textcircled{wrote} a letter.
    \item \underline{The cat} \textcircled{chased} the mouse.
\end{enumerate}
\end{quickcheck}

\begin{rememberbox}
You already know that every sentence has a subject and a verb (Chapter 1), and that the subject and verb must agree (Chapter 13). Now you will learn that the subject of a sentence does not always have to be the one \textit{doing} the action.
\end{rememberbox}

\section{Chapter Opener}
Ms. Sharma has written two versions of the same classroom news on the board.

\textbf{Version 1:} "Rohan broke the window."
\textbf{Version 2:} "The window was broken."

\textbf{Rohan:} \textit{(looking worried)} "Wait, both of those are about the same window! But the second one doesn't say who did it!"

\textbf{Maya:} "That's true... but it still makes sense, doesn't it? We know a window was broken, even without knowing who broke it."

\textbf{Ms. Sharma:} "Exactly, Maya! Sometimes we want to focus on \textit{who} does something, and sometimes we want to focus on \textit{what happens} to something — even if we don't know or don't want to say who did it. Let's explore these two ways of writing a sentence."

\section{Language Experience}
\textbf{Two Reports of the Same Event}

\begin{quote}
\textbf{Report A:} The gardener watered the plants every morning. He also cut the grass on Sundays.
\end{quote}

\begin{quote}
\textbf{Report B:} The plants were watered every morning. The grass was cut on Sundays.
\end{quote}

\section{Look and Notice}
\begin{thinkbox}
\begin{enumerate}
    \item In Report A, who is doing the watering and cutting?
    \item In Report B, can you still tell what happened to the plants and the grass?
    \item Does Report B tell you \textit{who} watered the plants?
    \item Which report would you use if you didn't know or didn't need to say who did the work?
\end{enumerate}
\end{thinkbox}

\section{Discover / Explicit Teaching}
Every action has two possible "starting points" for a sentence:

\begin{itemize}
    \item The \textbf{doer} — the person or thing that performs the action.
    \item The \textbf{receiver} — the person or thing that the action happens to.
\end{itemize}

\begin{table}[h!]
\centering
\begin{tabular}{|l|l|l|}
\hline
\textbf{Voice} & \textbf{Focus} & \textbf{Example} \\
\hline
\textbf{Active Voice} & The doer is the subject. & \textbf{Rohan} kicked the ball. \\
\textbf{Passive Voice} & The receiver is the subject. & The ball \textbf{was kicked} by Rohan. \\
\hline
\end{tabular}
\end{table}

\section{Grammar Word}
\begin{grammarword}{ACTIVE VOICE AND PASSIVE VOICE}
\textbf{Active voice} is when the subject of the sentence performs the action.
\textbf{Passive voice} is when the subject of the sentence receives the action.

\textit{Examples:}
\begin{itemize}
    \item Active: The chef \textbf{cooked} the meal.
    \item Passive: The meal \textbf{was cooked} by the chef.
\end{itemize}

\textit{Non-example:} "The meal cooked the chef" is not a correct passive of the first sentence — it reverses the actual doer and receiver, which changes the meaning entirely.

\textit{Quick Check:} In "The letter was posted by Maya," who is the receiver, and who is the doer? \textit{(Receiver: the letter. Doer: Maya.)}
\end{grammarword}

\section{Core Explanation}

\subsection*{Active voice — doer first}
``\textbf{The dog} chased the ball.'' \textit{Check:} Is the subject of the sentence the one doing the action? Then it is active.

\subsection*{Passive voice — receiver first}
``\textbf{The ball} was chased by the dog.'' \textit{Check:} Is the subject of the sentence the one the action happens to? Then it is passive.

\subsection*{Forming the passive (simple present and simple past only)}
Example (present): Maya \textbf{writes} the letter. $\rightarrow$ The letter \textbf{is written} by Maya.
Example (past): Maya \textbf{wrote} the letter. $\rightarrow$ The letter \textbf{was written} by Maya.
The passive always uses a form of "be" (is/am/are/was/were) plus the past participle of the main verb (written, cooked, broken, cleaned).

\section{Worked Example}
\begin{examplebox}
\textbf{Sentence to transform:} "The children ate the cake." (Change to passive voice.)

\textit{Step 1:} Identify the doer and the receiver. Doer: the children. Receiver: the cake.
\textit{Step 2:} Make the receiver the new subject: "The cake..."
\textit{Step 3:} Choose the correct form of "be" to match the tense (past) and the new subject (singular): "was."
\textit{Step 4:} Add the past participle of the main verb: "eat" $\rightarrow$ "eaten."
\textit{Step 5:} Add "by + doer" if needed: "by the children."

\textbf{Answer:} "\textbf{The cake was eaten by the children.}"
\end{examplebox}

\section{Guided Practice}
\begin{activitybox}{Activity A: Identification}
Say whether each sentence is active or passive.
\begin{enumerate}
    \item Rohan painted the fence.
    \item The fence was painted by Rohan.
    \item The cake was baked by Maya.
\end{enumerate}
\end{activitybox}

\begin{activitybox}{Activity B: Completion}
Change each active sentence into the passive voice.
\begin{enumerate}
    \item "The cat chased the mouse." $\rightarrow$ \_\_\_
    \item "Maya writes a story." $\rightarrow$ \_\_\_
\end{enumerate}
\end{activitybox}

\section{Independent Practice}
\begin{activitybox}{Independent Practice}
\begin{enumerate}
    \item Change this active sentence into the passive voice: "The teacher checked our homework."
    \item Change this passive sentence into the active voice: "The ball was thrown by Rohan."
    \item Write one sentence of your own in the active voice, then rewrite it in the passive voice.
    \item \textit{Reasoning:} Why might a school notice say "Homework must be submitted by Friday" instead of naming exactly who must submit it? What does the passive voice let the writer leave out?
\end{enumerate}
\end{activitybox}

\section{Common Confusion}
\begin{warningbox}
Learners often use the wrong form of the main verb in the passive (for example, "was eat" or "was breaking" instead of "was eaten" or "was broken"), because they forget that the passive needs the \textbf{past participle}, not the base form or the "-ing" form.

\textit{How to check:} Ask, "Have I used 'be' (is/am/are/was/were) followed by the special past-participle form of the verb (eaten, broken, written), not the plain verb or the '-ing' form?"
\end{warningbox}

\section{Summary}
\begin{summarybox}
\begin{itemize}
    \item \textbf{Active voice} focuses on the doer of an action; the doer is the subject.
    \item \textbf{Passive voice} focuses on the receiver of an action; the receiver is the subject.
    \item The passive voice is formed with a form of "be" (is/am/are/was/were) plus the past participle of the main verb.
    \item A "by" phrase can name the doer in a passive sentence, but it can also be left out.
    \item Passive voice is useful when the doer is unknown, unimportant, or already understood.
\end{itemize}
\end{summarybox}

\section{Key Terms}
\begin{table}[h!]
\centering
\begin{tabular}{|l|l|}
\hline
\textbf{Term} & \textbf{Simple Meaning} \\
\hline
Active voice & The subject performs the action \\
Passive voice & The subject receives the action \\
Doer & The person or thing performing the action \\
Receiver & The person or thing the action happens to \\
Past participle & The special verb form used after "be" in the passive (e.g., eaten, broken, written) \\
\hline
\end{tabular}
\end{table}

\begin{selfcheck}
I can:
\begin{itemize}
    \item identify the doer and receiver in a sentence.
    \item tell whether a sentence is active or passive.
    \item change a simple active sentence into the passive voice.
    \item change a simple passive sentence into the active voice.
    \item recognise passive voice in real-life writing.
\end{itemize}
\end{selfcheck}